% Cote 119: offprints, reports, courses, CV and letters (typed and printed part): archive cote 119, site page 267 (PDF page 268).
% typescript. Draft compared: none.
% AI-transcribed by Claude Opus 5.5 (claude-opus-5-5, adaptive thinking, effort high), Message Batches API, batch r119_full (msgbatch_01DwZNm5Y7YsKBWyAN4sFs8w), 2026-11-13; not yet checked by eye.
\scanpage{119/267}
\mspage{14}

trouve (dans l'optique où nous nous plaçons) que la proportion normale des
assistants et maîtres-assistants dans l'ensemble des enseignants du département
de Mathématique devrait être de l'ordre de un huitième. Il faut compter sans
\unsure{doute;compte|doute,compte} tenu des transitions nécessaires et de la dynamique de l'accrois-%
sement régulier des besoins en effectifs d'enseignants, que pendant assez
longtemps cette proportion resterait de l'ordre de un cinquième à un quart.
C'est à peu près la proportion inverse de celle qui est de règle à l'heure
actuelle ! Il est clair qu'un tel renversement n'est concevable que moyennant
une transformation radicale des méthodes d'enseignement actuellement en vigueur
à l'Université, où un petit état-major de professeurs et de maîtres de confé-%
rences dirige de haut une armée d'assistants et de maîtres-assistants dont
bien peu (tout comme dans l'armée moins pacifique qui fait la fierté de la
Nation) portent dans leur sac de fameux baton de Maréchal ! Une telle évolution
semble liée indissolublement à la disparition du système des cours magistraux,
disparition que les étudiants réclament avec insistance, mais que de nombreux
professeurs auront sans doute beaucoup de mal à admettre, et même à imaginer...

6. Dans l'immédiat, en plus des obstacles psychologiques sans doute extrê-%
mement sérieux de la part des professeurs, même parmi ceux qui sont le plus
sincèrement acquis au courant de renovation universitaire parti des étudiants,
il faut bien reconnaître qu'il y a également d'autres obstacles à une appli-%
cation conséquente de la nouvelle conception présentée ici du rôle de l'assis-%
tant et du mapitre-assistant. Il y a d'abord, bien entendu, la question des
crédits, un maître-assistant ``coûtant moins cher'' à l'Etat qu'un maître de
conférences. Les derniers évènements, il est vrai, permettent de se faire une
idée si les ``économies'' à court terme ainsi réalisées par l'Etat sont rentables
à longue échéance. Un problème sans doute plus délicat, car en grande partie
indépendant du bon vouloir de l'Etat dans les années à venir, est celui de
savoir combien parmi les actuels assistants et maîtres-assistants auraient un
niveau intellectuel suffisant pour pouvoir, dans les deux ou trois années qui
viennent, obtenir un doctorat d'enseignement et exercer les fonctions de
maître de conférences (mention enseignement). Il est probable que nombre d'entre
eux, (parmi ceux notamment qui s'épuisent actuellement en vains efforts pour

\mspage{\hand{267}}
